\section{The finite-section gap and its remainder classes}
\label{subsec:finite-gap-bridge}

This subsection isolates the passage from the two finite one-sided traces to
the whole-line pairings in
\cref{eq:shared-leading-gap,eq:shared-melnikov}.  It is useful to keep this
passage separate from the coefficient calculation: the latter determines the
number \(r_N/\alpha\), whereas the estimates below determine whether that
number is the derivative of the selected finite-section gap.

Write
\[
 c_H=\frac{2}{\alpha e},\qquad S=S_\delta,
 \qquad R=S+B,
\]
and let \(\widetilde Q_{N,\delta,\nu,\zeta}\) be the prepared planar field
which drives the selected traces on
\(J_a=[-R,0]\) and \(J_r=[0,R]\).  On
\(J_a^c=[-S,0]\) and \(J_r^c=[0,S]\), the traces and every history
backtrack entering the graph lie in the uncut region, so there
\(\widetilde Q=Q\) and
\cref{eq:shared-graph-expansion} applies.  Define the piecewise trace
\[
 z^{\rm pw}(s)=
 \begin{cases}
  z^a(s),&-R\le s\le0,\\
  z^r(s),&0\le s\le R,
 \end{cases}
\]
where the two values at \(s=0\) are kept distinct when an integral is split
there.  For \(\ell=1,2\), set
\begin{align*}
 A_{\ell,N}(z,\nu,\zeta)
 &=c_H\nabla\mathscr H_\alpha(z)^T
       q_{\ell,N}(z,\nu,\zeta),\\
 A_{3,N}(z,\delta,\nu,\zeta)
 &=c_H\nabla\mathscr H_\alpha(z)^T
       \mathcal R_{3,N}(z,\delta,\nu,\zeta).
\end{align*}
The unused parameter arguments are suppressed below; in particular,
\(q_{1,N}\) is independent of \(\zeta\).  Since
\(c_H\nabla\mathscr H_\alpha(\gamma_0(s))=\psi(s)\), introduce two
whole-line pairings and the finite baseline second-order pairing
\begin{align}
 M_{1,N}(\nu)
 &=\int_{\mathbb R}A_{1,N}(\gamma_0(s),\nu)\,\dd s,
                                                        \label{eq:gap-M1}\\
 B_{2,N}^{S}(\nu)
 &=\int_{-S}^{S}A_{2,N}(\gamma_0(s),\nu,0)\,\dd s,
                                                        \label{eq:gap-B2}\\
 M_{2,N}
 &=\int_{\mathbb R}\partial_\zeta
       A_{2,N}(\gamma_0(s),\nu,0)\,\dd s.             \label{eq:gap-M2}
\end{align}
The identities already obtained in
\cref{eq:shared-leading-gap,eq:shared-melnikov} say
\begin{equation}
 M_{1,N}(\nu)=\sqrt{2\pi}(\nu-\nu_{0,N}),
 \qquad
 M_{2,N}=\sqrt{2\pi}\frac{r_N}{\alpha}.              \label{eq:gap-leading-pairings}
\end{equation}
In particular, \(M_{2,N}\) is independent of \(\nu\).

We first record exactly which growing-tube estimates the gap calculation
uses.  Put \(W_{c,m}(s)=e^{c|s|}\langle s\rangle^m\), and let
\[
 \mathfrak A=\{\Id,\partial_\nu,\partial_\zeta,
 \partial_{\nu\nu},\partial_{\nu\zeta},
 \partial_{\zeta\zeta}\}
\]
denote the identity and the parameter derivatives needed below.

\begin{proposition}[Finite-section gap bridge]
\label[proposition]{prop:finite-gap-bridge}
Fix the common model bounds, the interval \(I_\nu\), and one fixed
admissible preparation with \(p>4\) satisfying
\cref{def:invdet-tame-preparation}.  By
\cref{prop:invdet-coefficient-envelope}, there are fixed
\(c,m,C_q\), independent of \(N\), for which the two graph coefficients
satisfy
\begin{equation}
 \left|D_z^a\partial_\nu^i\partial_\zeta^j
 q_{\ell,N}(\gamma_0(s)+\xi,\nu,\zeta)\right|
 \le C_q W_{c,m}(s),                                  \label{eq:gap-coefficient-envelope}
\end{equation}
for \(\ell=1,2\), \(0\le a\le3\),
\(0\le i,j\le2\), whenever
\(|s|\le S+1\), \(\nu\in I_\nu\), \(|\zeta|\le\zeta_*\), and
\(|\xi|\le C\delta W_{c,m}(s)\).  The same estimate for the derivatives of
\(\mathcal R_{3,N}\) listed in \cref{eq:shared-graph-remainder} is
\cref{eq:invdet-local-remainder}.  Then there are
\(\delta_b,\zeta_b,C_b>0\), depending only on the fixed common data and the
displayed envelope but not on \(N\), such that, for every \(N\ge2\),
\(0<\delta<\delta_b\), \(\nu\in I_\nu\), and
\(|\zeta|<\zeta_b\), the exact decomposition
\begin{align}
 \mathscr D_N^{\rm can}
 ={}&\delta M_{1,N}(\nu)+\delta^2B_{2,N}^{S}(\nu)
       +\delta^2\zeta M_{2,N}\notag\\
 &+E_{\partial,N}+E_{\rm sing,N}+E_{\rm mov,N}
       +E_{\rm par,N}+E_{\rm jet,N}+E_{\rm prep,N}       \label{eq:gap-exact-decomposition}
\end{align}
holds, where the six terms on the second line are given explicitly in
\cref{eq:gap-endpoint-error,eq:gap-singular-error,eq:gap-moving-error,%
eq:gap-parameter-error,eq:gap-jet-error,eq:gap-preparation-error} below.
Their parameter derivatives obey
\begin{align}
 E_{\partial,N}&=0,                                    \label{eq:gap-error-endpoint-bound}\\
 |\mathscr D E_{\rm sing,N}|+|\mathscr D E_{\rm prep,N}|
   &\le C_b\delta^4,
       &&\mathscr D\in\mathfrak A,                    \label{eq:gap-error-tail-bound}\\
 |E_{\rm mov,N}|+|\partial_\nu E_{\rm mov,N}|
       +|\partial_{\nu\nu}E_{\rm mov,N}|
   &\le C_b\delta^2,                                  \label{eq:gap-error-moving-nu}\\
 |\partial_\zeta E_{\rm mov,N}|
       +|\partial_{\nu\zeta}E_{\rm mov,N}|
       +|\partial_{\zeta\zeta}E_{\rm mov,N}|
   &\le C_b\delta^3,                                  \label{eq:gap-error-moving-zeta}\\
 |\mathscr D E_{\rm jet,N}|&\le C_b\delta^3,
       &&\mathscr D\in\mathfrak A,                    \label{eq:gap-error-jet-bound}\\
 |E_{\rm par,N}|+|\partial_\nu E_{\rm par,N}|
       +|\partial_{\nu\nu}E_{\rm par,N}|
   &\le C_b\delta^2\zeta^2,                          \label{eq:gap-error-parameter-zero}\\
 |\partial_\zeta E_{\rm par,N}|
       +|\partial_{\nu\zeta}E_{\rm par,N}|
   &\le C_b\delta^2|\zeta|,
 \qquad
 |\partial_{\zeta\zeta}E_{\rm par,N}|
   \le C_b\delta^2.                                  \label{eq:gap-error-parameter-derivatives}
\end{align}
Consequently, uniformly in the same quantifier order,
\begin{align}
 \frac{\mathscr D_N^{\rm can}}{\delta}
 &=\sqrt{2\pi}(\nu-\nu_{0,N})+O(\delta),             \label{eq:gap-bridge-expansion}\\
 \partial_\nu\mathscr D_N^{\rm can}
 &=\delta\sqrt{2\pi}+O(\delta^2),                    \label{eq:gap-bridge-nu}\\
 \partial_\zeta\mathscr D_N^{\rm can}
 &=\delta^2\sqrt{2\pi}\frac{r_N}{\alpha}
       +O(\delta^3+\delta^2|\zeta|),                 \label{eq:gap-bridge-zeta}\\
 \partial_{\zeta\zeta}\mathscr D_N^{\rm can}
 &=O(\delta^2),                                       \label{eq:gap-bridge-zz}\\
 |\partial_{\nu\nu}\mathscr D_N^{\rm can}|
       +|\partial_{\nu\zeta}\mathscr D_N^{\rm can}|
 &\le C_b\delta^2.                                   \label{eq:gap-bridge-mixed}
\end{align}
These are precisely the five estimates used in the implicit-function
argument.
\end{proposition}

\begin{proof}
We first derive the finite-section identity without replacing either trace
by the singular orbit.  The function \(\mathscr H_\alpha\) is a first
integral of \(q_0\), hence
\(\nabla\mathscr H_\alpha(z)^Tq_0(z)=0\) for every planar state \(z\).
The fundamental theorem of calculus on the two one-sided traces gives
\begin{align}
 \mathscr D_N^{\rm can}
 ={}&c_H\{\mathscr H_\alpha(z^a(-R))
              -\mathscr H_\alpha(z^r(R))\}\notag\\
 &+c_H\int_{-R}^0
     \nabla\mathscr H_\alpha(z^a(s))^T
       \widetilde Q(z^a(s))\,\dd s
  +c_H\int_0^R
     \nabla\mathscr H_\alpha(z^r(s))^T
       \widetilde Q(z^r(s))\,\dd s.                 \label{eq:gap-finite-section-identity}
\end{align}
Thus both outer endpoints and both one-sided traces occur in the exact
identity.  There are no moving integration limits when \(\nu\) or
\(\zeta\) is differentiated: \(R=S_\delta+B\) depends only on the fixed
value of \(\delta\).  Moreover, the selected boundary equations hold for
every \((\nu,\zeta)\), so
\begin{equation}
 E_{\partial,N}:=c_H\{\mathscr H_\alpha(z^a(-R))
              -\mathscr H_\alpha(z^r(R))\}=0          \label{eq:gap-endpoint-error}
\end{equation}
together with all of its \(\nu\)- and \(\zeta\)-derivatives.  This is exact
endpoint cancellation, not asymptotic endpoint suppression.
No derivative with respect to \(\delta\), for which the endpoint locations
would move, is asserted in this proposition.

On \([-S,S]\), insert
\cref{eq:shared-graph-expansion} into
\cref{eq:gap-finite-section-identity}.  Add and subtract the corresponding
singular-orbit integrands, extend the latter to \(\mathbb R\), and Taylor
expand only the \(\zeta\)-dependence of \(A_{2,N}\) at the fixed singular
orbit.  This gives \cref{eq:gap-exact-decomposition}, with
\begin{align}
 E_{\rm sing,N}
 ={}&-\delta\int_{|s|>S}A_{1,N}(\gamma_0(s),\nu)\,\dd s\notag\\
 &-\delta^2\zeta\int_{|s|>S}
       \partial_\zeta A_{2,N}(\gamma_0(s),\nu,0)\,\dd s,
                                                        \label{eq:gap-singular-error}\\
 E_{\rm mov,N}
 ={}&\delta\int_{-S}^{S}\!{}^{\rm pw}
       \{A_{1,N}(z^{\rm pw}(s),\nu)
               -A_{1,N}(\gamma_0(s),\nu)\}\,\dd s\notag\\
 &+\delta^2\int_{-S}^{S}\!{}^{\rm pw}
       \{A_{2,N}(z^{\rm pw}(s),\nu,\zeta)
               -A_{2,N}(\gamma_0(s),\nu,\zeta)\}\,\dd s,
                                                        \label{eq:gap-moving-error}\\
 E_{\rm par,N}
 ={}&\delta^2\int_{-S}^{S}
       \bigl\{A_{2,N}(\gamma_0(s),\nu,\zeta)
             -A_{2,N}(\gamma_0(s),\nu,0)\notag\\
 &\hspace{42mm}
             -\zeta\partial_\zeta
                    A_{2,N}(\gamma_0(s),\nu,0)\bigr\}\,\dd s,
                                                        \label{eq:gap-parameter-error}\\
 E_{\rm jet,N}
 ={}&\delta^3\int_{-S}^{S}\!{}^{\rm pw}
       A_{3,N}(z^{\rm pw}(s),\delta,\nu,\zeta)\,\dd s,
                                                        \label{eq:gap-jet-error}\\
 E_{\rm prep,N}
 ={}&c_H\int_{[-R,-S]\cup[S,R]}\!{}^{\rm pw}
       \nabla\mathscr H_\alpha(z^{\rm pw}(s))^T
       \{\widetilde Q(z^{\rm pw}(s))-q_0(z^{\rm pw}(s))\}\,\dd s.
                                                        \label{eq:gap-preparation-error}
\end{align}
Here \(\int^{\rm pw}\) means that the negative half uses \(z^a\) and the
positive half uses \(z^r\).  The quantities in
\cref{eq:gap-singular-error,eq:gap-moving-error,%
eq:gap-parameter-error,eq:gap-jet-error,eq:gap-preparation-error} are
definitions, so no remainder has yet been estimated or hidden.

We now prove the stated bounds.  The trace estimates in
\cref{eq:shared-trace-jets} imply the pointwise estimates
\begin{align}
 |z^{\rm pw}(s)-\gamma_0(s)|
       +|\partial_\nu z^{\rm pw}(s)|
 &\le C\delta W_{c,m}(s),\notag\\
 |\partial_{\nu\nu}z^{\rm pw}(s)|
       +|\partial_\zeta z^{\rm pw}(s)|
       +|\partial_{\nu\zeta}z^{\rm pw}(s)|
       +|\partial_{\zeta\zeta}z^{\rm pw}(s)|
 &\le C\delta^2 W_{c,m}(s).                          \label{eq:gap-pointwise-trace-jets}
\end{align}
Increasing \(c,m\) finitely many times is harmless.  It also gives, for
example,
\begin{equation}
 \int_{-R}^{R}e^{-s^2/2}
       |\partial_\zeta z^{\rm pw}(s)|^2
       \langle s\rangle^{m}\,\dd s\le C\delta^4,     \label{eq:gap-quadratic-trace}
\end{equation}
because
\(\int_{\mathbb R}e^{-s^2/2+2c|s|}
\langle s\rangle^{3m}\,\dd s<\infty\).

For sufficiently small \(\delta\),
\(\sup_{|s|\le R}\delta W_{c,m}(s)=o(1)\).  Direct differentiation of
\cref{eq:shared-first-integral} therefore gives, for every fixed
\(0\le k\le4\),
\begin{equation}
 |D^k\mathscr H_\alpha(\gamma_0(s)+\xi)|
 \le C e^{-s^2/2+c|s|}\langle s\rangle^m
 \quad\text{if }|\xi|\le C\delta W_{c,m}(s).          \label{eq:gap-H-envelope}
\end{equation}
Indeed every such derivative is \(e^{-2\alpha Y}\) times a fixed
polynomial in \((X,Y)\), while
\(Y=(s^2-2)/(4\alpha)+o(1)\) on this tube.  Combining
\cref{eq:gap-coefficient-envelope,eq:shared-graph-remainder,%
eq:gap-pointwise-trace-jets,eq:gap-H-envelope} with the mean-value theorem
and the chain rule gives the following orders after integration:
\begin{center}
\begin{tabular}{c@{\qquad}cccccc}
\toprule
 & \(1\) & \(\partial_\nu\) & \(\partial_\zeta\)
 & \(\partial_{\nu\nu}\) & \(\partial_{\nu\zeta}\)
 & \(\partial_{\zeta\zeta}\)\\
\midrule
\(E_{\rm mov,N}\)
 & \(\delta^2\) & \(\delta^2\) & \(\delta^3\)
 & \(\delta^2\) & \(\delta^3\) & \(\delta^3\)\\
\(E_{\rm jet,N}\)
 & \(\delta^3\) & \(\delta^3\) & \(\delta^3\)
 & \(\delta^3\) & \(\delta^3\) & \(\delta^3\)\\
\bottomrule
\end{tabular}
\end{center}
For example, the potentially leading structural derivative of the first
moving-trace integral is
\[
 \delta\int_{-S}^{S}\!{}^{\rm pw}
 D_zA_{1,N}(z^{\rm pw}(s),\nu)
       \partial_\zeta z^{\rm pw}(s)\,\dd s,
\]
because \(\partial_\zeta q_{1,N}=0\).  It is \(O(\delta^3)\), not
\(O(\delta^2)\), by
\cref{eq:gap-pointwise-trace-jets}.  A second \(\zeta\)-derivative adds
either \(\partial_{\zeta\zeta}z^{\rm pw}=O_G(\delta^2)\) or the quadratic
term controlled by \cref{eq:gap-quadratic-trace}.  The \(\nu\nu\)-column
instead contains
\(\delta D_z\partial_\nu A_{1,N}\partial_\nu z^{\rm pw}\), which is only
\(O(\delta^2)\).  These observations give
\cref{eq:gap-error-moving-nu,eq:gap-error-moving-zeta}; applying the same
chain rule to \(A_{3,N}\) gives \cref{eq:gap-error-jet-bound}.

Taylor's formula with integral remainder yields the pointwise identity
\begin{align*}
 &A_{2,N}(\gamma_0(s),\nu,\zeta)
       -A_{2,N}(\gamma_0(s),\nu,0)
       -\zeta\partial_\zeta A_{2,N}(\gamma_0(s),\nu,0)\\
 &\qquad
 =\zeta^2\int_0^1(1-t)
       \partial_{\zeta\zeta}A_{2,N}
          (\gamma_0(s),\nu,t\zeta)\,\dd t.
\end{align*}
The Gaussian envelope and its first two \(\nu\)-derivatives are integrable
on \(\mathbb R\).  Differentiating this formula once or twice in \(\zeta\)
proves
\cref{eq:gap-error-parameter-zero,eq:gap-error-parameter-derivatives}.

It remains to estimate the two tail classes.  For every fixed \(c,m\),
\begin{equation}
 \int_{|s|>S}e^{-s^2/2+c|s|}\langle s\rangle^m\,\dd s
 \le C e^{-S^2/2+cS}\langle S\rangle^{m+1}
 =\delta^{p-o(1)}.                                   \label{eq:gap-gaussian-tail-bound}
\end{equation}
This proves the asserted estimate for \(E_{\rm sing,N}\).  On each
preparation interval, \(S\le|s|\le S+B\), the interval length is fixed.
The preparation derivative bound in
\cref{def:canonical-preparation}, the trace bounds
\cref{eq:gap-pointwise-trace-jets}, and
\cref{eq:gap-H-envelope} show, after any
\(\mathscr D\in\mathfrak A\), that the integrand in
\cref{eq:gap-preparation-error} is bounded by
\(C e^{-s^2/2+c'|s|}\langle s\rangle^{m'}\).  Thus the same calculation
applies to \(E_{\rm prep,N}\).  Since \(p>4\), decrease \(\delta_b\) so
that both quantities are at most \(C_b\delta^4\).  This proves
\cref{eq:gap-error-tail-bound}.

Finally, \cref{eq:gap-coefficient-envelope,eq:gap-H-envelope} make
\(B_{2,N}^{S}\) and its first two \(\nu\)-derivatives uniformly bounded,
independently of the receding value of \(S\).
Insert \cref{eq:gap-leading-pairings} and
\cref{eq:gap-error-endpoint-bound,eq:gap-error-tail-bound,%
eq:gap-error-moving-nu,eq:gap-error-moving-zeta,%
eq:gap-error-jet-bound,eq:gap-error-parameter-zero,%
eq:gap-error-parameter-derivatives} into
\cref{eq:gap-exact-decomposition}.  The five conclusions
\cref{eq:gap-bridge-expansion,eq:gap-bridge-nu,eq:gap-bridge-zeta,%
eq:gap-bridge-zz,eq:gap-bridge-mixed} follow term by term.
\end{proof}

\begin{remark}[Audit of the bridge inputs]
\label[remark]{rem:gap-bridge-audit}
The fixed admissible preparation provides the analogues of the
usual selected-trace assumptions as follows.
The hitting-time statement is replaced by the definition of the fixed
intervals \([-R,0]\) and \([0,R]\).  The strong trace bounds
\cref{eq:shared-trace-jets} imply all Gaussian \(L^1\) bounds in
\cref{eq:gap-pointwise-trace-jets,eq:gap-quadratic-trace}.  Finally, the
endpoint equation \(\mathscr H_\alpha=0\) is an identity for the selected
trace family, so its parameter derivatives vanish exactly; no separate
endpoint-growth estimate is needed.

The only non-trace input specific to this bridge is the pointwise
growing-tube bound \cref{eq:gap-coefficient-envelope} for \(q_{1,N}\),
\(q_{2,N}\), and their state and mixed parameter derivatives.  It is exactly
\cref{prop:invdet-coefficient-envelope}; the corresponding graph-remainder
bound is \cref{eq:invdet-local-remainder}.  Hence the detailed graph
construction and the trace estimates together supply every hypothesis of
\cref{prop:finite-gap-bridge}.  No additional physical-outer trace or
ambient RFDE inverse is required for this fixed-preparation finite-section
bridge.
\end{remark}
